\documentclass{article}
\usepackage[utf8]{inputenc}
\usepackage[russian]{babel}
\usepackage{amsmath, amssymb, amsthm}

\theoremstyle{definition}
\newtheorem{definition}{Определение}
\newtheorem{theorem}{Теорема}
\newtheorem{lemma}{Лемма}
\newtheorem{corollary}{Следствие}

\begin{document}

\title{Конспект по теме: \\ Функции нескольких переменных, \\ частные производные, градиент, \\ экстремумы, метод множителей Лагранжа, \\ полные системы функций}
\author{По учебнику В.Ф. Бутузова \\ «Лекции по математическому анализу. Части II и III»}
\date{}
\maketitle

\section*{Введение}
Конспект составлен на основе полного текста учебника В.Ф. Бутузова (2014 г.). Освещены разделы, соответствующие вопросам задания: функция многих переменных, частные производные, градиент, дифференциал, задачи безусловной и условной оптимизации, локальный экстремум, метод множителей Лагранжа, функция Лагранжа, а также (на основе части III) полные системы функций, разложение по полной системе и остаточный член.

\section{Функция нескольких переменных (основные понятия)}
\subsection{$m$-мерное евклидово пространство}
Множество упорядоченных совокупностей $m$ чисел $(x_1,\dots,x_m)$ называется $m$-мерным координатным пространством. Расстояние между точками $M_1(x_1,\dots,x_m)$ и $M_2(y_1,\dots,y_m)$ определяется формулой
\[
\rho(M_1,M_2)=\sqrt{(x_1-y_1)^2+\dots+(x_m-y_m)^2}.
\]
Пространство с таким расстоянием называется $m$-мерным евклидовым пространством $\mathbb{R}^m$.

\subsection{Функция многих переменных}
Если каждой точке $M(x_1,\dots,x_m)$ из множества $\{M\}\subset\mathbb{R}^m$ поставлено в соответствие число $u$, то говорят, что на $\{M\}$ определена функция $m$ переменных. Обозначения: $u=f(x_1,\dots,x_m)$ или $u=f(M)$.
Для функции двух переменных $z=f(x,y)$ графиком является поверхность в пространстве $Oxyz$.

\subsection{Предел и непрерывность}
Определения предела по Коши и по Гейне, критерий Коши, непрерывность по совокупности переменных и по отдельной переменной. Непрерывность по совокупности влечёт непрерывность по каждой переменной, но обратное неверно (пример $u=\frac{xy}{x^2+y^2}$ с доопределением нулём в нуле).

\subsection{Основные теоремы о непрерывных функциях}
\textbf{Теорема 7 (арифметические операции).} Если $f(M)$ и $g(M)$ непрерывны в точке $A$, то $f\pm g$, $fg$, $f/g$ (при $g(A)\neq0$) непрерывны в $A$.

\textbf{Теорема 8 (непрерывность сложной функции).} Композиция непрерывных функций непрерывна.

\textbf{Теорема 9 (устойчивость знака).} Если $f$ непрерывна в $A$ и $f(A)>0$, то существует окрестность, в которой $f(M)>0$.

\textbf{Теорема 10 (о прохождении через промежуточное значение).} Непрерывная на связном множестве функция принимает все промежуточные значения между значениями в любых двух точках.

\textbf{Лемма 3.} Если $\{M_n\}\subset\{M\}$ и $\{M_n\}\to A$, а множество $\{M\}$ замкнуто, то $A\in\{M\}$.

\textbf{Теорема 11 (первая теорема Вейерштрасса).} Непрерывная на замкнутом ограниченном множестве функция ограничена.
\begin{proof}
Предположим противное: $f$ не ограничена. Тогда $\forall n$ существует $M_n$ такой, что $|f(M_n)|>n$. Последовательность $\{M_n\}$ ограничена, поэтому из неё можно выделить сходящуюся подпоследовательность $M_{k_n}\to A$. По лемме 3 $A$ принадлежит множеству. В силу непрерывности $f(M_{k_n})\to f(A)$, но $|f(M_{k_n})|>k_n\to\infty$, противоречие. \end{proof}

\textbf{Теорема 12 (вторая теорема Вейерштрасса).} Непрерывная на замкнутом ограниченном множестве функция достигает своих точных верхней и нижней граней.
\begin{proof}
Докажем для верхней грани $U=\sup f$. По определению верхней грани существует последовательность $M_n$ такая, что $f(M_n)\to U$. Из ограниченности последовательности $M_n$ выделяем сходящуюся подпоследовательность $M_{k_n}\to A$. В силу замкнутости $A$ принадлежит множеству. По непрерывности $f(A)=\lim f(M_{k_n})=U$, так что $U$ достигается. \end{proof}

\textbf{Теорема 13 (Кантора).} Непрерывная на замкнутом ограниченном множестве функция равномерно непрерывна.
Доказательство проводится методом от противного с использованием теоремы Больцано–Вейерштрасса (см. учебник).

\section{Частные производные и дифференцируемость}
\subsection{Частные производные}
Частная производная функции $u=f(x_1,\dots,x_m)$ в точке $M$ по переменной $x_k$ определяется как
\[
\frac{\partial u}{\partial x_k}(M)=\lim_{\Delta x_k\to0}\frac{\Delta_{x_k}u}{\Delta x_k},
\]
где $\Delta_{x_k}u=f(x_1,\dots,x_k+\Delta x_k,\dots,x_m)-f(x_1,\dots,x_k,\dots,x_m)$.
\textbf{Физический смысл:} скорость изменения функции в направлении оси $Ox_k$.
Существование частных производных не гарантирует непрерывности (пример: функция, равная 1 на осях координат и 0 в остальных точках).

\subsection{Дифференцируемость функции многих переменных}
Функция $u=f(x_1,\dots,x_m)$ называется \textbf{дифференцируемой} в точке $M$, если её полное приращение можно представить в виде
\[
\Delta u = A_1\Delta x_1+\dots+A_m\Delta x_m + \alpha_1\Delta x_1+\dots+\alpha_m\Delta x_m,
\]
где $A_i$ — числа, не зависящие от $\Delta x_i$, а $\alpha_i\to0$ при $\Delta x_1\to0,\dots,\Delta x_m\to0$ и $\alpha_i(0,\dots,0)=0$. Это условие эквивалентно
\[
\Delta u = A_1\Delta x_1+\dots+A_m\Delta x_m + o(\rho),\quad \rho=\sqrt{(\Delta x_1)^2+\dots+(\Delta x_m)^2}.
\]
Из дифференцируемости следует непрерывность.

\textbf{Теорема 14 (необходимое условие).} Если функция дифференцируема в точке $M$, то она имеет в этой точке частные производные по всем переменным, причём $A_k=\frac{\partial u}{\partial x_k}(M)$.
Обратное неверно.

\begin{proof}[Доказательство теоремы 14]
Запишем условие дифференцируемости в виде (9.5):
\[
\Delta u = A_1\Delta x_1 + \dots + A_m\Delta x_m + \alpha_1\Delta x_1 + \dots + \alpha_m\Delta x_m.
\]
Положим все $\Delta x_i = 0$, кроме $\Delta x_k$, и $\Delta x_k \neq 0$. Тогда
\[
\Delta u = \Delta_{x_k}u = A_k\Delta x_k + \alpha_k\Delta x_k,
\]
где $A_k$ — число, а $\alpha_k \to 0$ при $\Delta x_k \to 0$. Следовательно,
\[
\frac{\Delta_{x_k}u}{\Delta x_k} = A_k + \alpha_k \to A_k \quad \text{при } \Delta x_k \to 0.
\]
Таким образом, существует предел $\lim_{\Delta x_k\to0}\frac{\Delta_{x_k}u}{\Delta x_k} = A_k$, то есть существует частная производная $\frac{\partial u}{\partial x_k}(M) = A_k$ для любого $k=1,\dots,m$. \end{proof}

\textbf{Теорема 15 (достаточное условие).} Если функция имеет частные производные в некоторой окрестности точки $M$ и эти производные непрерывны в самой точке $M$, то функция дифференцируема в $M$.

\begin{proof}[Доказательство теоремы 15 (для функции двух переменных)]
Возьмём $\Delta x$, $\Delta y$ столь малыми, чтобы точка $M_1(x+\Delta x,y+\Delta y)$ лежала в $\varepsilon$-окрестности точки $M(x,y)$. Полное приращение запишем в виде
\[
\Delta u = f(x+\Delta x,y+\Delta y)-f(x,y) = \bigl[f(x+\Delta x,y+\Delta y)-f(x,y+\Delta y)\bigr] + \bigl[f(x,y+\Delta y)-f(x,y)\bigr].
\]
К каждой разности применим формулу Лагранжа:
\[
\Delta u = f'_x(x+\theta_1\Delta x, y+\Delta y)\,\Delta x + f'_y(x, y+\theta_2\Delta y)\,\Delta y,\quad 0<\theta_1,\theta_2<1.
\]
В силу непрерывности частных производных в точке $M$ имеем
\[
f'_x(x+\theta_1\Delta x, y+\Delta y)=f'_x(x,y)+\alpha_1,\quad f'_y(x, y+\theta_2\Delta y)=f'_y(x,y)+\alpha_2,
\]
где $\alpha_1,\alpha_2\to0$ при $\Delta x,\Delta y\to0$. Подставляя, получаем
\[
\Delta u = f'_x(x,y)\Delta x + f'_y(x,y)\Delta y + \alpha_1\Delta x + \alpha_2\Delta y,
\]
что и означает дифференцируемость. \end{proof}

\subsection{Дифференциал функции многих переменных}
\textbf{Дифференциалом} (первым дифференциалом) функции $u=f(M)$ в точке $M$ называется линейная часть приращения:
\[
du = \frac{\partial u}{\partial x_1}(M)dx_1+\dots+\frac{\partial u}{\partial x_m}(M)dx_m,
\]
где $dx_i=\Delta x_i$ — дифференциалы независимых переменных.

\textbf{Лемма об инвариантности формы первого дифференциала:} формула для $du$ сохраняет свой вид и в том случае, когда $x_1,\dots,x_m$ являются дифференцируемыми функциями других переменных (однако теперь $dx_i$ — дифференциалы этих функций, и вообще $dx_i\neq\Delta x_i$).

Правила дифференцирования для дифференциалов аналогичны правилам для функций одной переменной.

\subsection{Геометрический смысл дифференцируемости. Касательная плоскость}
Для функции двух переменных $z=f(x,y)$ дифференцируемость в точке $M_0(x_0,y_0)$ означает существование касательной плоскости к графику функции в точке $N_0(x_0,y_0,z_0)$, $z_0=f(x_0,y_0)$. Уравнение касательной плоскости:
\[
Z-z_0 = \frac{\partial z}{\partial x}(M_0)(x-x_0) + \frac{\partial z}{\partial y}(M_0)(y-y_0).
\]
Вектор нормали: $\vec{n}=\left\{\frac{\partial z}{\partial x}(M_0),\frac{\partial z}{\partial y}(M_0),-1\right\}$.

\begin{proof}[Доказательство теоремы 17]
Пусть $N(x,y,z)\in S$, $z=f(x,y)$. Обозначим $\Delta x=x-x_0$, $\Delta y=y-y_0$, $\Delta z=z-z_0$. Из дифференцируемости
\[
\Delta z = A\Delta x + B\Delta y + o(\rho),\quad A=\frac{\partial z}{\partial x}(M_0),\;B=\frac{\partial z}{\partial y}(M_0),\;\rho=\sqrt{\Delta x^2+\Delta y^2}.
\]
Рассмотрим плоскость $P$: $Z-z_0 = A(x-x_0)+B(y-y_0)$. Для точки $N_2(x,y,Z)$ на этой плоскости имеем $\rho(N,N_2)=|z-Z|=o(\rho)$. Так как $\rho(N,N_1)\le\rho(N,N_2)$ (перпендикуляр короче наклонной), а $\rho(N,N_0)\ge\rho$, то
\[
\frac{\rho(N,N_1)}{\rho(N,N_0)}\le\frac{o(\rho)}{\rho}\to0,
\]
что доказывает, что $P$ — касательная плоскость. \end{proof}

\section{Производная по направлению и градиент}
\subsection{Производная по направлению}
Пусть функция $u=f(x,y,z)$ определена в окрестности точки $M_0$, $\vec{l}=(\cos\alpha,\cos\beta,\cos\gamma)$ — единичный вектор, задающий направление. Производная по направлению $\vec{l}$ в точке $M_0$ определяется как
\[
\frac{\partial u}{\partial l}(M_0)=\lim_{t\to0}\frac{f(M_0+t\vec{l})-f(M_0)}{t}.
\]
Если функция дифференцируема, то
\[
\frac{\partial u}{\partial l}(M_0)=\frac{\partial u}{\partial x}(M_0)\cos\alpha+\frac{\partial u}{\partial y}(M_0)\cos\beta+\frac{\partial u}{\partial z}(M_0)\cos\gamma. \tag{9.15}
\]

\begin{proof}[Вывод формулы (9.15)]
На прямой $M_0+t\vec{l}$ имеем $x=x_0+t\cos\alpha$, $y=y_0+t\cos\beta$, $z=z_0+t\cos\gamma$. Сложная функция $\varphi(t)=f(x_0+t\cos\alpha,y_0+t\cos\beta,z_0+t\cos\gamma)$ дифференцируема при $t=0$, и
\[
\varphi'(0)=\frac{\partial f}{\partial x}(M_0)\cos\alpha+\frac{\partial f}{\partial y}(M_0)\cos\beta+\frac{\partial f}{\partial z}(M_0)\cos\gamma.
\]
Но $\varphi'(0)=\lim_{t\to0}\frac{\varphi(t)-\varphi(0)}{t}=\frac{\partial u}{\partial l}(M_0)$. \end{proof}

\subsection{Градиент}
\textbf{Градиентом} дифференцируемой функции $u=f(x,y,z)$ в точке $M_0$ называется вектор
\[
\operatorname{grad} u(M_0)=\frac{\partial u}{\partial x}(M_0)\vec{i}+\frac{\partial u}{\partial y}(M_0)\vec{j}+\frac{\partial u}{\partial z}(M_0)\vec{k}.
\]
Тогда производная по направлению выражается как скалярное произведение:
\[
\frac{\partial u}{\partial l}(M_0)=\bigl(\operatorname{grad}u(M_0)\cdot\vec{l}\bigr).
\]
Отсюда следует:
\begin{itemize}
    \item $\frac{\partial u}{\partial l}(M_0)$ максимальна, когда $\vec{l}$ совпадает по направлению с $\operatorname{grad}u(M_0)$, и равна $|\operatorname{grad}u(M_0)|$.
    \item Вектор градиента указывает направление наискорейшего роста функции в данной точке, а его модуль — скорость роста в этом направлении.
    \item Градиент ортогонален поверхности уровня (линии уровня) функции.
\end{itemize}

\subsection{Примеры физических полей}
Потенциал электростатического поля $u(M)=k e/r$, напряжённость $\vec{E}=-\operatorname{grad}u$. Ньютоновский потенциал поля тяготения $u(M)=\gamma m/r$, сила $\vec{F}=\operatorname{grad}u$.

\section{Частные производные и дифференциалы высших порядков}
Если частная производная $\frac{\partial u}{\partial x_i}$ дифференцируема, то её частные производные называются производными второго порядка:
\[
\frac{\partial^2 u}{\partial x_k\partial x_i}=\frac{\partial}{\partial x_k}\left(\frac{\partial u}{\partial x_i}\right).
\]
При $k\neq i$ — смешанные производные.

\textbf{Теорема 18 (о равенстве смешанных производных).} Если смешанные производные $f_{xy}''$ и $f_{yx}''$ определены в окрестности точки $M_0$ и непрерывны в этой точке, то они равны: $f_{xy}''(M_0)=f_{yx}''(M_0)$. Более сильная форма: если функция дважды дифференцируема в точке $M_0$, то смешанные производные второго порядка не зависят от порядка дифференцирования.

\begin{proof}[Доказательство теоремы 18 (для функции двух переменных)]
Рассмотрим квадрат со стороной $h$ с вершиной в $M_0(x_0,y_0)$. Введём функцию
\[
F(h)=f(x_0+h,y_0+h)-f(x_0+h,y_0)-f(x_0,y_0+h)+f(x_0,y_0).
\]
С одной стороны, $F(h)=\varphi(x_0+h)-\varphi(x_0)$, где $\varphi(x)=f(x,y_0+h)-f(x,y_0)$. Применяя дважды формулу Лагранжа, получаем
\[
F(h)=f_{xy}''(x_0+\theta_1h,y_0+\theta_2h)\,h^2.
\]
С другой стороны, $F(h)=\psi(y_0+h)-\psi(y_0)$ с $\psi(y)=f(x_0+h,y)-f(x_0,y)$, откуда
\[
F(h)=f_{yx}''(x_0+\theta_1'h,y_0+\theta_2'h)\,h^2.
\]
Переходя к пределу при $h\to0$ и используя непрерывность смешанных производных, получаем $f_{xy}''(x_0,y_0)=f_{yx}''(x_0,y_0)$. \end{proof}

\textbf{Дифференциалы высших порядков} определяются индуктивно: $d^2u=d(du)$ (при этом $dx,dy$ считаются постоянными при повторном дифференцировании). Для функции двух переменных
\[
d^2u = \frac{\partial^2 u}{\partial x^2}dx^2 + 2\frac{\partial^2 u}{\partial x\partial y}dxdy + \frac{\partial^2 u}{\partial y^2}dy^2.
\]
В общем случае $d^n u = \left(\frac{\partial}{\partial x_1}dx_1+\dots+\frac{\partial}{\partial x_m}dx_m\right)^n u$.

\section{Формула Тейлора}
Если функция $f(x_1,\dots,x_m)$ $(n+1)$ раз дифференцируема в окрестности точки $M_0$, то для любой точки $M$ из этой окрестности справедлива формула Тейлора:
\[
f(M)-f(M_0)=df\big|_{M_0}+\frac{1}{2!}d^2f\big|_{M_0}+\dots+\frac{1}{n!}d^nf\big|_{M_0}+\frac{1}{(n+1)!}d^{n+1}f\big|_{N},
\]
где $N$ — некоторая точка на отрезке $M_0M$, а дифференциалы вычисляются по формуле
\[
d^k f = \left(\frac{\partial}{\partial x_1}\Delta x_1+\dots+\frac{\partial}{\partial x_m}\Delta x_m\right)^k f.
\]
Для дважды дифференцируемой функции справедлива формула с остаточным членом в форме Пеано:
\[
\Delta f = df|_{M_0} + \frac{1}{2}d^2f|_{M_0} + o(\rho^2),\quad \rho=\rho(M_0,M).
\]

\begin{proof}[Набросок доказательства теоремы 19]
Вводим параметр $t$ вдоль отрезка $M_0M$: $x_i=x_i^0+t\Delta x_i$, $0\le t\le1$. Функция $F(t)=f(x_1^0+t\Delta x_1,\dots,x_m^0+t\Delta x_m)$ является функцией одной переменной $t$, $(n+1)$ раз дифференцируемой. Применяем к ней формулу Тейлора с остаточным членом в форме Лагранжа при $t_0=0$, $t=1$, $dt=1$:
\[
F(1)-F(0)=dF(0)+\frac{1}{2!}d^2F(0)+\dots+\frac{1}{n!}d^nF(0)+\frac{1}{(n+1)!}d^{n+1}F(\theta),\quad 0<\theta<1.
\]
Так как $x_i$ линейны по $t$, то $d^kF(0)=d^kf(M_0)$ и $d^{k+1}F(\theta)=d^{k+1}f(N)$, где $N$ соответствует $t=\theta$. \end{proof}

\section{Задача безусловной оптимизации (локальный экстремум)}
\subsection{Постановка задачи}
Пусть $f(M)$ — функция, определённая в некоторой области. Требуется найти точки, в которых $f$ принимает наибольшее или наименьшее значение в некоторой окрестности. Такие точки называются точками локального экстремума (максимума или минимума). Формально: $M_0$ — точка локального максимума, если существует окрестность $U(M_0)$, такая, что $\forall M\in U(M_0): f(M)\le f(M_0)$ (для минимума $\ge$).

\subsection{Необходимое условие (теорема 20)}
Если $f$ дифференцируема в точке $M_0$ и $M_0$ — точка локального экстремума, то все частные производные в этой точке равны нулю, т.е. $df|_{M_0}=0$. Доказательство приведено ранее.

\subsection{Достаточные условия (теоремы 21, 22)}
Для дважды дифференцируемой функции в стационарной точке $M_0$ исследуется знак второго дифференциала $d^2f|_{M_0}$. Если $d^2f|_{M_0}$ положительно определён — минимум, отрицательно определён — максимум, знакопеременен — экстремума нет (седловая точка). Доказательства приведены выше.

\section{Неявные функции}
\subsection{Одно уравнение}
Рассмотрим уравнение $F(x,y)=0$. Говорят, что оно определяет неявную функцию $y=f(x)$ на интервале $X$, если для любого $x\in X$ уравнение имеет единственное решение $y=f(x)$ и $F(x,f(x))\equiv0$.

\textbf{Теорема 2 (существование и непрерывность неявной функции).} Пусть
\begin{enumerate}
    \item $F(x,y)$ определена и непрерывна в окрестности точки $M_0(x_0,y_0)$;
    \item $F(x_0,y_0)=0$, $F_y'(x_0,y_0)\neq0$;
    \item частная производная $F_y'(x,y)$ непрерывна в точке $M_0$.
\end{enumerate}
Тогда существуют числа $d>0$, $c>0$ и единственная функция $y=f(x)$, определённая на интервале $(x_0-d,x_0+d)$, такая, что $f(x_0)=y_0$, $|f(x)-y_0|\le c$ и $F(x,f(x))\equiv0$. Функция $f(x)$ непрерывна на $(x_0-d,x_0+d)$.

\begin{proof}[Доказательство теоремы 2]
Для определённости считаем $F_y'(x_0,y_0)>0$. В силу непрерывности $F_y'$ найдётся прямоугольник $Q=\{ (x,y):|x-x_0|\le\tilde d, |y-y_0|\le c\}$, где $F_y'>0$, т.е. $F(x,y)$ строго возрастает по $y$. Так как $F(x_0,y_0)=0$, то $F(x_0,y_0-c)<0$, $F(x_0,y_0+c)>0$. В силу непрерывности по $x$ существует $d\le\tilde d$, такое, что $F(x,y_0-c)<0$, $F(x,y_0+c)>0$ при $|x-x_0|<d$. Тогда для каждого фиксированного $x$ из этого интервала существует единственное $y\in(y_0-c,y_0+c)$, для которого $F(x,y)=0$ (теорема о промежуточном значении и монотонность). Это и есть $y=f(x)$. Непрерывность $f(x)$ следует из устойчивости знака: для любого $\varepsilon>0$ выбираем $\delta$ так, чтобы знаки $F(x,f(x_0)\pm\varepsilon)$ оставались теми же, что и при $x=x_0$; тогда $f(x)$ попадает в $\varepsilon$-окрестность $f(x_0)$. \end{proof}

\textbf{Теорема 3 (дифференцируемость неявной функции).} Если в дополнение к условиям теоремы 2 функция $F(x,y)$ дифференцируема в точке $M_0$, то неявная функция $y=f(x)$ дифференцируема в точке $x_0$, и
\[
f'(x_0)=-\frac{F_x'(x_0,y_0)}{F_y'(x_0,y_0)}.
\]

\begin{proof}[Доказательство теоремы 3]
Из тождества $F(x,f(x))\equiv0$ и дифференцируемости $F$ в точке $M_0$ имеем
\[
0=F(x_0+\Delta x, f(x_0+\Delta x))-F(x_0,f(x_0))=F_x'(x_0,y_0)\Delta x+F_y'(x_0,y_0)\Delta f+\alpha\Delta x+\beta\Delta f,
\]
где $\alpha,\beta\to0$ при $\Delta x\to0$ и $\Delta f=f(x_0+\Delta x)-f(x_0)$. Отсюда
\[
\frac{\Delta f}{\Delta x}=-\frac{F_x'(x_0,y_0)+\alpha}{F_y'(x_0,y_0)+\beta}.
\]
Переходя к пределу при $\Delta x\to0$ и используя непрерывность $f$, получаем искомую формулу. \end{proof}

Аналогичные теоремы формулируются для неявной функции $y=f(x_1,\dots,x_n)$, заданной уравнением $F(x_1,\dots,x_n,y)=0$, и для систем неявных функций. В последнем случае ключевую роль играет якобиан $\frac{D(F_1,\dots,F_m)}{D(y_1,\dots,y_m)}$.

\section{Условный экстремум. Метод множителей Лагранжа}
\subsection{Постановка задачи}
Требуется найти экстремумы функции $u=f(x_1,\dots,x_n)$ при наличии $m$ условий связи ($m<n$):
\[
F_1(x_1,\dots,x_n)=0,\quad\dots,\quad F_m(x_1,\dots,x_n)=0.
\]
Точка $M_0$, удовлетворяющая условиям связи, называется точкой \textbf{условного экстремума}, если для всех точек $M$ из некоторой окрестности $M_0$, также удовлетворяющих условиям связи, выполняется соответствующее неравенство.

\subsection{Метод множителей Лагранжа}
Вводится \textbf{функция Лагранжа}:
\[
\Phi(x_1,\dots,x_n,\lambda_1,\dots,\lambda_m)=f(x_1,\dots,x_n)+\lambda_1F_1(x_1,\dots,x_n)+\dots+\lambda_mF_m(x_1,\dots,x_n).
\]
Числа $\lambda_1,\dots,\lambda_m$ называются \textbf{множителями Лагранжа}.

\textbf{Теорема 8 (необходимое условие условного экстремума).} Пусть выполнены условия:
\begin{enumerate}
    \item функции $f,F_1,\dots,F_m$ дифференцируемы в точке $M_0$;
    \item $F_j(M_0)=0$ ($j=1,\dots,m$);
    \item якобиан $\frac{D(F_1,\dots,F_m)}{D(x_1,\dots,x_m)}\big|_{M_0}\neq0$ (т.е. условия связи регулярны).
\end{enumerate}
Если $M_0$ — точка условного экстремума, то существуют числа $\lambda_1,\dots,\lambda_m$ такие, что
\[
\frac{\partial\Phi}{\partial x_i}(M_0)=0,\quad i=1,\dots,n.
\]

\begin{proof}[Доказательство теоремы 8 (по Лагранжу)]
В силу условия регулярности система $F_j=0$ определяет в окрестности $M_0$ единственную систему неявных функций
\[
x_1=\varphi_1(x_{m+1},\dots,x_n),\;\dots,\;x_m=\varphi_m(x_{m+1},\dots,x_n).
\]
Подставляя их в $f$, получаем функцию $g(x_{m+1},\dots,x_n)=f(\varphi_1,\dots,\varphi_m,x_{m+1},\dots,x_n)$. Точка $M_0$ соответствует $M_0'(x_{m+1}^0,\dots,x_n^0)$, и $g$ имеет в ней безусловный экстремум. Следовательно,
\[
dg|_{M_0'}=0.
\]
Выразим $dg$ через дифференциалы функции Лагранжа. Заметим, что в точках, удовлетворяющих условиям связи, $\Phi=f$. Поэтому
\[
dg = d\Phi = \sum_{i=1}^m\frac{\partial\Phi}{\partial x_i}dx_i+\sum_{i=m+1}^n\frac{\partial\Phi}{\partial x_i}dx_i.
\]
Выберем множители $\lambda_j$ так, чтобы коэффициенты при $dx_1,\dots,dx_m$ обратились в нуль в точке $M_0$:
\[
\frac{\partial f}{\partial x_i}(M_0)+\sum_{j=1}^m\lambda_j\frac{\partial F_j}{\partial x_i}(M_0)=0,\quad i=1,\dots,m.
\]
Это линейная система относительно $\lambda_j$; её определитель $\frac{D(F_1,\dots,F_m)}{D(x_1,\dots,x_m)}\big|_{M_0}\neq0$, поэтому $\lambda_j$ существуют и единственны. При таком выборе
\[
dg|_{M_0'}= \sum_{i=m+1}^n\frac{\partial\Phi}{\partial x_i}(M_0)dx_i.
\]
Так как $dx_{m+1},\dots,dx_n$ — дифференциалы независимых переменных, из $dg=0$ следует
\[
\frac{\partial\Phi}{\partial x_i}(M_0)=0,\quad i=m+1,\dots,n.
\]
Таким образом, все частные производные $\Phi$ по $x_i$ в точке $M_0$ равны нулю. \end{proof}

\subsection{Интерпретация множителей Лагранжа}
Множители Лагранжа можно интерпретировать как «цены» ограничений. При малом изменении правой части условия связи (замена $F_j=0$ на $F_j=\varepsilon$) изменение значения функции в точке условного экстремума пропорционально $\lambda_j$: $\Delta f\approx -\lambda_j\varepsilon$ (знак зависит от формулировки). В физических и экономических задачах $\lambda_j$ характеризуют чувствительность оптимума к изменению ограничений.

\subsection{Достаточные условия условного экстремума}
Для проверки, действительно ли найденная стационарная точка функции Лагранжа даёт условный экстремум, исследуют второй дифференциал функции Лагранжа при условии, что дифференциалы переменных связаны линейными соотношениями, полученными дифференцированием условий связи:
\[
\sum_{i=1}^n \frac{\partial F_j}{\partial x_i}(M_0)\,dx_i=0,\quad j=1,\dots,m.
\]
Если на этом множестве ненулевых наборов $dx_i$ квадратичная форма $d^2\Phi|_{M_0}$ является знакоопределённой (положительной — минимум, отрицательной — максимум), то точка $M_0$ даёт условный экстремум. Если форма знакопеременна — экстремума нет (седловая точка в рамках условной задачи).

\section{Седловые точки}
В теории оптимизации и вариационного исчисления седловая точка функции (в частности, функции Лагранжа) — это точка, в которой функция имеет минимум по одним переменным и максимум по другим. Для задачи условной оптимизации точка $(x^0,\lambda^0)$ называется седловой точкой функции Лагранжа, если
\[
\Phi(x^0,\lambda)\le\Phi(x^0,\lambda^0)\le\Phi(x,\lambda^0)
\]
для всех $x$ из допустимой области и $\lambda\ge0$ (для задач с ограничениями-неравенствами). В классической задаче с ограничениями-равенствами седловая точка функции Лагранжа (без ограничений на знак $\lambda$) соответствует точке условного экстремума.

\section{Полные системы функций. Разложение по полной системе. Остаточный член}
\subsection{Пространство $Q[a,b]$ и скалярное произведение}
Рассмотрим множество всех кусочно-непрерывных на сегменте $[a,b]$ функций, причём в точках разрыва значение функции полагается равным полусумме левого и правого пределов. Это множество образует линейное пространство (обозначим его $Q[a,b]$). Введём скалярное произведение:
\[
(f,g)=\int_a^b f(x)g(x)\,dx.
\]
Норма элемента определяется как $\|f\|=\sqrt{(f,f)}$.

\subsection{Ортогональные и ортонормированные системы}
Система элементов $\{\psi_n\}$ называется \textbf{ортогональной}, если $(\psi_i,\psi_j)=0$ при $i\ne j$. Если дополнительно $\|\psi_n\|=1$, система называется \textbf{ортонормированной}.

\textbf{Примеры:}
\begin{itemize}
    \item Тригонометрическая система $1,\cos x,\sin x,\cos 2x,\sin 2x,\dots$ ортогональна на $[-\pi,\pi]$.
    \item Система полиномов Лежандра $P_n(x)$ ортогональна на $[-1,1]$.
    \item Система полиномов Чебышёва $T_n(x)$ ортогональна на $[-1,1]$ с весом $1/\sqrt{1-x^2}$.
\end{itemize}

\subsection{Ряд Фурье по ортогональной системе}
Пусть $\{\psi_n\}$ — ортогональная система, не содержащая нулевых элементов. Для любого элемента $f\in Q[a,b]$ определим его \textbf{коэффициенты Фурье}:
\[
f_n = \frac{(f,\psi_n)}{\|\psi_n\|^2},\qquad n=1,2,\dots
\]
Ряд $\sum_{n=1}^{\infty} f_n\psi_n$ называется \textbf{рядом Фурье} элемента $f$ по системе $\{\psi_n\}$.

\subsection{Экстремальное свойство частичных сумм ряда Фурье}
\textbf{Теорема 2 (лекция 19).} При фиксированном $n$ среди всех линейных комбинаций $\sum_{k=1}^n c_k\psi_k$ наименьшее отклонение от $f$ по норме даёт частичная сумма ряда Фурье $S_n=\sum_{k=1}^n f_k\psi_k$.

\begin{proof}
Вычислим квадрат нормы разности:
\[
\Bigl\|\sum_{k=1}^n c_k\psi_k - f\Bigr\|^2 = \sum_{k=1}^n c_k^2\|\psi_k\|^2 - 2\sum_{k=1}^n c_k (f,\psi_k) + \|f\|^2.
\]
Подставим $(f,\psi_k)=f_k\|\psi_k\|^2$ и выделим полный квадрат:
\[
= \sum_{k=1}^n (c_k - f_k)^2\|\psi_k\|^2 + \|f\|^2 - \sum_{k=1}^n f_k^2\|\psi_k\|^2.
\]
Это выражение минимально, когда $c_k=f_k$ для всех $k$. \end{proof}

\subsection{Тождество и неравенство Бесселя}
Из доказательства при $c_k=f_k$ получаем \textbf{тождество Бесселя}:
\[
\|f\|^2 - \sum_{k=1}^n f_k^2\|\psi_k\|^2 = \|S_n - f\|^2 \ge 0.
\]
Отсюда следует \textbf{неравенство Бесселя}:
\[
\sum_{k=1}^n f_k^2\|\psi_k\|^2 \le \|f\|^2 \quad \forall n,
\]
а значит, ряд $\sum_{k=1}^{\infty} f_k^2\|\psi_k\|^2$ сходится и его сумма не превосходит $\|f\|^2$.

\subsection{Замкнутые и полные системы}
\begin{definition}
Ортогональная система $\{\psi_n\}$ называется \textbf{замкнутой}, если для любого элемента $f\in Q[a,b]$ и любого $\varepsilon>0$ существует линейная комбинация конечного числа элементов системы, отличающаяся от $f$ по норме менее чем на $\varepsilon$.
\end{definition}
В силу экстремального свойства это эквивалентно тому, что частичные суммы ряда Фурье сходятся к $f$ по норме: $\|S_n - f\|\to0$.

\begin{definition}
Ортогональная система $\{\psi_n\}$ называется \textbf{полной}, если единственным элементом, ортогональным всем $\psi_n$, является нулевой элемент:
\[
\forall n:\ (f,\psi_n)=0 \quad\Rightarrow\quad f\equiv0.
\]
\end{definition}

\textbf{Теорема 3 (критерий замкнутости для ортонормированной системы).} 
Для ортонормированной системы $\{\psi_n\}$ следующие условия эквивалентны:
\begin{enumerate}
    \item Система замкнута.
    \item Для любого элемента $f$ выполняется \textbf{равенство Парсеваля}:
    \[
    \sum_{n=1}^{\infty} f_n^2 = \|f\|^2,\qquad f_n=(f,\psi_n).
    \]
\end{enumerate}
\begin{proof}
Из тождества Бесселя $\|S_n-f\|^2 = \|f\|^2 - \sum_{k=1}^n f_k^2$. Сходимость $S_n$ к $f$ по норме равносильна стремлению правой части к нулю, что и есть равенство Парсеваля. Замкнутость означает возможность приближения, а для ортонормированной системы наилучшее приближение дают именно $S_n$, поэтому замкнутость эквивалентна сходимости $S_n$ к $f$, т.е. равенству Парсеваля. \end{proof}

\textbf{Теорема 4.} Любая замкнутая система является полной.
\begin{proof}
Пусть система $\{\psi_n\}$ замкнута и элемент $f$ ортогонален всем $\psi_n$. Тогда все коэффициенты Фурье $f_n=0$. Из равенства Парсеваля (которое выполняется в силу замкнутости) получаем $\|f\|^2=0$, следовательно $f\equiv0$. \end{proof}

\textbf{Теорема 5.} Если система $\{\psi_n\}$ полна, то два различных элемента не могут иметь одинаковые ряды Фурье.
\begin{proof}
Если $f$ и $g$ имеют одинаковые коэффициенты, то $f-g$ ортогонален всем $\psi_n$, и в силу полноты $f-g\equiv0$, т.е. $f=g$. \end{proof}

\subsection{Примеры полных систем}
\subsubsection{Тригонометрическая система}
Тригонометрическая система на $[-\pi,\pi]$ является замкнутой (и, следовательно, полной). Это следует из теоремы Вейерштрасса о равномерной аппроксимации непрерывных функций тригонометрическими многочленами и свойства сходимости в среднем. Для любой кусочно-непрерывной функции $f(x)$ её тригонометрический ряд Фурье сходится к $f$ в среднем, и выполняется равенство Парсеваля:
\[
\frac{a_0^2}{2} + \sum_{n=1}^{\infty} (a_n^2 + b_n^2) = \frac{1}{\pi}\int_{-\pi}^{\pi} f^2(x)\,dx,
\]
где $a_n,b_n$ — обычные коэффициенты Фурье.

\subsubsection{Алгебраические полиномы}
\textbf{Теорема Вейерштрасса (о приближении многочленами).} Любую непрерывную на сегменте $[a,b]$ функцию можно равномерно приблизить алгебраическим многочленом. Это означает, что система степеней $\{1,x,x^2,\dots\}$ полна в пространстве непрерывных функций (а значит, и в $Q[a,b]$) по норме равномерной сходимости, что влечёт её полноту и по норме $L^2$. Однако эта система не ортогональна; ортогональные полиномы (Лежандра, Чебышёва и др.) образуют полные ортонормированные системы на соответствующих отрезках.

\subsection{Разложение произвольной функции по полной системе}
Если $\{\psi_n\}$ — полная ортонормированная система в пространстве $Q[a,b]$, то любой элемент $f$ из этого пространства (в частности, любая кусочно-непрерывная функция) может быть разложен в ряд Фурье, сходящийся к $f$ в среднем:
\[
f(x) = \sum_{n=1}^{\infty} f_n \psi_n(x),\quad f_n = (f,\psi_n).
\]
Сходимость в среднем означает:
\[
\lim_{n\to\infty} \int_a^b \bigl(f(x) - S_n(x)\bigr)^2 dx = 0,
\]
где $S_n(x) = \sum_{k=1}^n f_k\psi_k(x)$.

\subsection{Остаточный член разложения}
Остаточный член $R_n$ при замене $f$ частичной суммой $S_n$ в смысле сходимости в среднем равен
\[
R_n = \|f - S_n\|^2 = \|f\|^2 - \sum_{k=1}^n f_k^2.
\]
В силу равенства Парсеваля $R_n\to0$ при $n\to\infty$. Для тригонометрического ряда это означает, что среднеквадратичное отклонение частичной суммы от функции стремится к нулю. Для конкретных функций часто можно получить и более точные оценки, например, используя гладкость функции, можно оценить скорость убывания коэффициентов Фурье и, следовательно, остаточный член.

\end{document}
